\documentclass[pdflatex,compress,9pt,
	xcolor={dvipsnames,dvipsnames,svgnames,x11names,table},
	hyperref={colorlinks = true,breaklinks = true, urlcolor = NavyBlue, breaklinks = true}]{beamer}
	
\usepackage[utf8]{inputenc}
\usepackage[T2A,T1]{fontenc}
%\usefonttheme{structuresmallcapsserif}
\usepackage{bookman}

\usetheme{Warsaw}
%\usecolortheme{seahorse}Teal DarkCyan
%\usecolortheme{wolverine} 
\usecolortheme[named=Crimson]{structure}

%%%%%%%%%%%%%%%%%%%%%%%%%%%%%%%%%%%

\usepackage{gensymb} % degree symbol
\usepackage[super]{nth}
%%%%%%%%%%%%%%%%%%%%%%%%%%%%%%%%%%%
% ----------------------------------------------------------------------------
% *** START BIBLIOGRAPHY <<<
% ----------------------------------------------------------------------------
\usepackage[
	backend=biber, 
%	style = numeric,
	style=apa,
	maxbibnames=99,
%	citestyle=authoryear,
	citestyle=numeric,
	giveninits=true,
	isbn=true,
	url=true,
	natbib=true,
	sorting=ndymdt,
	bibencoding=utf8,
	useprefix=false,
	language=auto, 
	autolang=other,
	backref=true,
	backrefstyle=none,
	indexing=cite,
]{biblatex}
\DeclareSortingTemplate{ndymdt}{
  \sort{
    \field{presort}
  }
  \sort[final]{
    \field{sortkey}
  }
  \sort{
    \field{sortname}
    \field{author}
    \field{editor}
    \field{translator}
    \field{sorttitle}
    \field{title}
  }
  \sort[direction=descending]{
    \field{sortyear}
    \field{year}
    \literal{9999}
  }
  \sort[direction=descending]{
    \field[padside=left,padwidth=2,padchar=0]{month}
    \literal{99}
  }
  \sort[direction=descending]{
    \field[padside=left,padwidth=2,padchar=0]{day}
    \literal{99}
  }
  \sort{
    \field{sorttitle}
  }
  \sort[direction=descending]{
    \field[padside=left,padwidth=4,padchar=0]{volume}
    \literal{9999}
  }
}

\addbibresource{Minsk.bib}
\renewcommand*{\bibfont}{\tiny} % \tiny \scriptsize \footnotesize

\setbeamertemplate{bibliography item}{\insertbiblabel}

% ----------------------------------------------------------------------------
% *** END BIBLIOGRAPHY <<<
% ----------------------------------------------------------------------------
% Путь к файлам с иллюстрациями
\graphicspath{{fig/}} % path to folder with Figures

\title{Biblioteca Salaborsa \\
Bologna, Italia}\vspace{2em}
 
\subtitle{\vspace*{0.5cm}Presentato a: \\
231 Corso italiano semestrale Livello: B1.3\\
Insegnante: Losi Liliana \\
Piazza San Giovanni in Monte 4\\
Bologna, Ilatia}

\author{Polina Lemenkova}
\date{dicembre 10, 2024}

\begin{document}
\begin{frame}
           \titlepage
\end{frame}

\section{Introduction}

\subsection{Study Area}
\begin{frame}{Study Area: P\"{a}rnu Region, Estonia}
\begin{minipage}[0.4\textheight]{\textwidth}
\begin{columns}[T]
\begin{column}{0.5\textwidth}
\begin{figure}[H]
	\centering
		\includegraphics[width=4.2cm]{Fig_01.jpg}
\end{figure}
\begin{block}{Study Area}
The research region encompasses coastal area of Baltic Sea: south-western Estonia
\end{block}
 
\begin{alertblock}{Spatial extent}
Spatial extent of the study area is limited to the surroundings of P\"{a}rnu County.
\end{alertblock}

\end{column}
\begin{column}{0.5\textwidth}
\begin{figure}[H]
	\centering
		\includegraphics[width=4.5cm]{Fig_01.jpg}
\end{figure}
\end{column}
\end{columns}
\end{minipage}
\end{frame}

\subsection{Landscapes}
\begin{frame}\frametitle{Landscapes}

\begin{block}{South-West Estonia: Unique Environment}
South-west Estonia is known for unique environmental settings: mild maritime climate, broad beaches, coniferous pine forests on the coastal zone
\end{block}

\begin{alertblock}{Landscapes}
Landscapes here are rich and world-known for their diversity, variability, unique composition structure and high esthetic value
\end{alertblock}

\begin{examples}{Types of Landscapes}
Landscape types include, for example, mixed and broadleaved forests, traditional agricultural semi-natural landscapes, wooded meadows, plant communities, heathland, bogs and moors, complex anthropogenic areas with different land use structure, shrubland, grasslands, birch-dominating coastal areas and flooded meadows
\end{examples}

\end{frame}

\section{Examples of Landscapes}
\subsection{Baltic Sea Coasts}
\begin{frame}\frametitle{Examples of Landscapes: Baltic Sea Coasts}
\begin{examples}{Landscapes:}
Baltic Sea Coasts. Photos: author.
\end{examples}
\begin{figure}[H]
	\centering
		\includegraphics[width=9.0cm]{Fig_01.jpg}
\end{figure}
\end{frame}

\subsection{Marine Landscapes: P\"{a}rnu}
\begin{frame}\frametitle{Marine Landscapes: P\"{a}rnu Coasts}
\begin{examples}{Marine Landscapes:}
 P\"{a}rnu Coasts. Photos: author.
\end{examples}
\begin{figure}[H]
	\centering
		\includegraphics[width=9.0cm]{Fig_01.jpg}
\end{figure}
\end{frame}

\subsection{Lacustrine landscapes: Luitemaaa}
\begin{frame}\frametitle{Lacustrine landscapes: Luitemaa Nature Conservation Area}
\begin{examples}{Lacustrine landscapes:}
Luitemaa Nature Conservation Area. Photos: author.
\end{examples}
\begin{figure}[H]
	\centering
		\includegraphics[width=9.0cm]{Fig_01.jpg}
\end{figure}
\end{frame}

\section{Methods}
\begin{frame}\frametitle{Methods}

\begin{block}{GIS Projects}
The GIS projects has been organized and executed in two different software: \alert{Arc GIS 10.0} and \alert{IDRISI GIS Andes 15.0}.
\end{block}

\begin{block}{Landsat TM}
The raster processing GIS approach and classification was applied in the current work towards Landsat TM two images."
\end{block}

\begin{block}{ISOCLUST}
The method is based on the ISOCLUST unsupervised classification executed by means of IDRISI GIS.
\end{block}

\begin{block}{Machine Learning}
The ISOCLUST available in IDRISI GIS performs the most of the image processing workflow in semi-automatically regime
\end{block}

\begin{block}{Land Cover Classes}
It results in a map with pre-defined number of 16 land class categories which enable to compare two different stages of landscape development: \alert{“earlier”} and \alert{“now”}.
\end{block}

\begin{block}{Objectivity}
The ISOCLUST method was chosen, since it enables to avoid subjectivity in classification.
\end{block}

\end{frame}

\begin{frame}\frametitle{Workflow Step-1}

\begin{block}{IMPORT}
Initially, the data of Landsat TM were imported to IDRISI Andes GIS from \alert{GeoTIFF} format to IDRISI specific format \alert{.rst}, through \alert{Data Provider Format} import.\\ As each Landsat TM scene is a multispectral image with several spectral bands, each band was displayed and visualized as a separate image.
\end{block}

\begin{block}{COLOR COMPOSITE}
Afterwords, the images were composed using \alert{Color Composite} function. \\
The combination of three bands was made as a single color composite image (bands 2-3-4).\\
 This composition displays urban areas distinctively, which enables to clearly recognize them
\end{block}

\begin{block}{PROJECT}
Then data were organized in a created project in \alert{IDRISI GIS}.
\end{block}

\end{frame}

\begin{frame}\frametitle{Land Cover Classes}
\begin{minipage}[0.4\textheight]{\textwidth}
\begin{columns}[T]
\begin{column}{0.5\textwidth}
\small{\begin{itemize}
	\item discontinuous urban fabric
	\item industrial or commercial units
	\item green urban areas
	\item pastures
	\item complex cultivation patterns
	\item agriculture lands with grass
	\item broad-leaved forest
	\item coniferous forest
	\item mixed forest
	\item natural grassland
	\item moors and heathlands
	\item transitional woodland
	\item beaches, dunes, sand
	\item island marshes
	\item water bodies
	\item sea and ocean
\end{itemize}}
\end{column}
\begin{column}{0.5\textwidth}
\begin{figure}[H]
	\centering
		\includegraphics[width=4.5cm]{Fig_01.jpg}
\end{figure}
\end{column}
\end{columns}
\end{minipage}
\end{frame}

\section{Thanks}
\begin{frame}{Thanks}
\emph{Thank you for attention !}\\
\begin{figure}[H]
	\centering
		\includegraphics[width=8.0cm]{Fig_01.jpg}
\end{figure}

\normalsize
Acknowledgement: \\
The research has been done at the University of Tartu \\
under support of DoRa grant (ESF, Estonia).\\
The University of Tartu provided \alert{data} and \alert{software}:\\
 \alert{CORINE} vector layers, \alert{IDRISI GIS} 15.0 and \alert{ArcGIS} 10.0.
\end{frame}
	
%%%%%%%%%%% Bibliography %%%%%%%	

\end{document}
%Changing the font size locally (from biggest to smallest):	
%\Huge
%\huge
%\LARGE
%\Large
%\large
%\normalsize (default)
%\small
%\footnotesize
%\scriptsize
%\tiny

\end{document}